\chapter{我们为什么开始做这轮优化}

这份文档记录的是为了跑通并加速综合编译负载，对 WaterOS 内核进行的一轮实际优化。优化对象不是某一个孤立函数，而是一条贯穿用户程序、系统调用、任务调度、虚拟内存、VFS、文件系统和块设备的长路径。单次调用中的一次查表、一次页表刷新或一次内存清零并不起眼，但 BuildStorm 会重复创建进程、装载 ELF、查询工具链文件并申请匿名页，这些固定成本最终会累计成数百秒的差距。

仓库中的 \code{perf/*} 分支保留了这轮工作的实验轨迹。分支既包括地址空间缓存、条件 TLB 刷新、只读页共享和缓存准入等有效方案，也包括扩大缓存、修改调度周期、重写 LRU 或增加专用小对象池等失败实验。我们没有把“做过”直接等同于“有效”：正文以已经进入主线、或具有完整 A/B 证据的方案为主；未获得端到端收益的实验仍保留，用来解释为什么没有继续沿那个方向投入。

整个过程可以概括为三件事。先把运行环境和完成标记固定下来，保证两次结果确实可比；再通过 \code{pc-hot}、\code{wait-hot} 和 syscall profile 从完整负载中提取指令、等待和参数分布；最后只改变一个主要因素做交错 A/B，并在功能回归通过后决定保留还是回退。这样做比直接改数据结构慢一些，却能区分“符号自身变冷”和“用户真正更快”这两件事。

\section{优化问题的分层}

本轮工作逐步收敛为四层问题。第一层是重复工作，例如同一路径被多次解析、同一 fd 在一条系统调用链中反复进入 registry。第二层是缓存策略，例如工具链只读页的跨进程复用，以及低命中块缓存是否值得接纳第一次访问。第三层是同步粒度，例如 lazy VMA 没有产生 PTE 时仍进行全地址空间 TLB 刷新。第四层是工作迁移，即把必须完成、但不必由繁忙 CPU 同步完成的页面清零转移到 idle CPU。

这些层次之间有先后关系。若调用链仍在重复查找，扩大缓存只会掩盖上层浪费；若缓存对象缺少稳定身份，跨进程共享就无法保证失效语义；若没有先测量 PTE 是否变化，简单删除 fence 又会破坏正确性。因此，后续章节按“测量---消除重复---改善复用---缩小同步---迁移前台工作”的顺序展开，而不是按提交日期罗列。

\section{实验分支提供了什么}

\begin{table}[H]
  \centering
  \caption{性能实验分支的主要类别}
  \begin{tabularx}{0.96\textwidth}{p{3.1cm}p{4.2cm}X}
    \toprule
    类别 & 代表性方向 & 在本文中的用途 \\
    \midrule
    测量工具 & syscall profile、cache diagnostics、ELF page profile & 先确认调用次数、参数形态、命中率和复用距离，再选择优化对象。 \\
    已保留优化 & 当前地址空间缓存、单页 COW 刷新、mprotect 变化摘要、只读 mmap/ELF 页缓存 & 构成地址空间、TLB 和文件页复用的主要实现路径。 \\
    策略实验 & block-cache admission、page-cache capacity、negative dentry、readahead & 用完整 A/B 确定容量、准入和淘汰策略，而不是依赖经验参数。 \\
    回退实验 & 40/48 MiB 页缓存、128 B TLSF 固定池、调度周期修改、intrusive LRU & 记录无收益、退化或正确性问题，防止重复尝试同一路径。 \\
    \bottomrule
  \end{tabularx}
\end{table}

\section{最终纳入的优化范围}

\begin{table}[H]
  \centering
  \caption{正文覆盖的最终优化主线}
  \begin{tabularx}{0.96\textwidth}{p{3.2cm}X}
    \toprule
    方向 & 最终保留的工作 \\
    \midrule
    测量与基线 & 固定 QEMU 9.2.1、镜像、CPU 数量和 \code{-snapshot}，建立完整 BuildStorm 与短窗口诊断两级方法。 \\
    syscall 与 FD & 批量导入用户元数据，复用已分类 socket 和 FD I/O lease，减少用户地址空间及 registry 的重复访问。 \\
    VFS 与文件系统 & 复用路径解析结果，缓存正负 dentry，改进缓存准入，并为稳定只读访问建立明确契约。 \\
    exec 与内存映射 & 共享只读 ELF/可执行 mmap 物理页，复用 exec 文件句柄，并按真实 PTE 变化决定 TLB 同步。 \\
    任务与同步 & 保留当前地址空间快路径、单页 COW 失效和 futex registry 分片等能够证明语义边界的改动。 \\
    物理页分配 & 由 idle CPU 预先清零页面并放入零页池，前台 fault 和页表分配优先消费已清零页面。 \\
    \bottomrule
  \end{tabularx}
\end{table}
